\subsection{Hub tasksets and storyworlds with a local 27B model}
\label{sec:hub}
The SWE-bench run used a 4B model. To see how the packet behaves with a stronger local agent on other long tasks, we
ran Ternary Bonsai 2 27B for one night on the RTX 5080. The model is quantised to PTQ1\_0 (5.9 GB) and served by the
PrismML llama-server fork:
\begin{itemize}
\item two slots of 32,768 tokens each;
\item a thinking budget of 1,536 tokens per turn;
\item temperature 0.6.
\end{itemize}
The environments are verifiers packages from the Prime Intellect Environments Hub plus the MoralityLab storyworld
environments. Each ran unmodified in its own virtual environment, through a relay that records the server's exact
token counts.

Two relay arms differ only in what the model sees:
\begin{itemize}
\item \textbf{raw} forwards the history unchanged;
\item \textbf{goblin} applies the MCP-03 packet scorer to the growing history once it passes 6,000 estimated tokens. The
task and the last two assistant turns stay verbatim. Earlier turns become an action ledger plus the 1,800 tokens of
earlier lines that score highest under BM25 and the exact-match channel.
\end{itemize}
The two arms ran side by side on the same tasks. The learned routers are not used here, because no environment exposes
a code repository.

\paragraph{Storyworlds.}
Table~\ref{tab:hub-storyworlds} covers the held-out lanes; no training or confirmatory rows were run:
\begin{itemize}
\item \textbf{Jinn--Beast lanes.} Bonsai 27B clears all four:
  \begin{itemize}
  \item cross-frame decisions: mean reward \HubJbmCrossDevEightZeroMean{} on \HubJbmCrossDevEightZeroN{} tasks;
  \item construct and moral-reasoner lanes: \HubJbmConstructsDevMean{} and \HubJbmReasonerDevMean{};
  \item the tool-using control mesh: \HubJbmMeshTwoDevMean{} on \HubJbmMeshTwoDevN{} tasks, with the registered target
  action every time.
  \end{itemize}
  No lane produced a critical action.
\item \textbf{storyworld-constitution.} The model fails it, with reward \HubSwConstitutionEvalMean{} on
\HubSwConstitutionEvalN{} tasks:
  \begin{itemize}
  \item the prompt specifies the output as \texttt{<option\_id>}, and the model wrapped every answer in angle brackets,
  which the environment's strict parser rejects;
  \item with the brackets removed, its choices match the frontier option on \HubSwcBracket{} of tasks, close to chance
  over three options.
  \end{itemize}
\end{itemize}

\paragraph{Agentic tasksets.}
Table~\ref{tab:hub-agentic} gives the outcomes. Bonsai 27B solves:
\begin{itemize}
\item every hard-tier feishu-office task in both arms;
\item \HubNotebookHardenedSolved{} of \HubNotebookHardenedN{} durable-notebook episodes;
\item \HubSqliteAllSolved{} of \HubSqliteAllN{} mcp-sqlite questions, scored by the taskset's exact match because its
API judge was removed.
\end{itemize}
The packet result is mixed and mostly negative:
\begin{itemize}
\item \textbf{event-recon (tier 1).} The raw history overflowed the 32K context in all \HubErTOneRawN{} episodes, so the
environment ended each one with an error. The goblin arm never overflowed and ran every episode to the 40-turn cap,
but it never filed a report either. Instead it re-read files the packet had dropped. Both arms scored zero, and the
goblin arm used \HubErTOneRatio{} times the raw arm's prompt tokens because its episodes ran four times longer.
\item \textbf{doc-edit.} Here the history fits in context, and compaction hurt. The raw arm solved
\HubDoceditRawSolved{} of \HubDoceditRawN{} tasks, the goblin arm \HubDoceditGoblinSolved{}. Raw was better on
\HubDoceditRawBetter{} tasks and goblin on \HubDoceditGoblinBetter{}. The goblin arm re-read pages, left stale edit sites
and ran out of turns, using \HubDoceditRatio{} times the prompt tokens.
\item \textbf{feishu-office and mcp-sqlite.} Prompts stayed under the threshold, so the packet never engaged. The
feishu pairs (7/7 against 7/7 on the hard tier, and 7/7 against 6/7 on the expert tier) therefore measure sampling
variation only.
\item \textbf{Late compaction (event-recon, tier 1).} After the main queue, a third arm compacted only past 18,000
estimated tokens, with a 6,000-token packet and four kept turns. It overflowed in all \HubErTOneLateN{} episodes
anyway:
  \begin{itemize}
  \item the character estimate (3.2 characters per token) undercounts the number-heavy tool output, so real prompts
  passed 32,768 tokens before the estimate reached the threshold;
  \item four kept turns of large file reads were enough to overflow on their own.
  \end{itemize}
\end{itemize}
Compacting early trades away detail that a 27B agent still needs. On these tasks it does not save tokens, because the
agent spends more turns recovering that detail. Compacting late, without an exact token count, does not stop
overflows. A packet for an agent near its context limit needs the server's exact token count, and it should keep the
specific items the agent has already acted on, such as files read and sites edited, not only the lines that score
highest against the task.
